\ifdefined\gkver\else\def\gkver{student}\fi
\documentclass[\gkver]{gaokaozhenti}
\begin{document}

\chapter{2006年广东}

\begin{enumerate}
\item
%% number: 1
%% paperName: 2006年广东高考第 1 题 4 分
%% typeId: 1
%% type: 单选题
%% chapter: 运动和力
%% point: 牛顿第一定律
%% method: 无
%% score: 4
%% degree: 650
%% duplicateId: 0
%% body:
下列对运动的认识不正确的是\xzanswer{A}

\fourchoices[answer=A]
{亚里士多德认为物体的自然状态是静止的，只有当它受到力的作用才会运动}
{伽利略认为力不是维持物体速度的原因}
{牛顿认为力的真正效应总是改变物体的速度，而不仅仅是使之运动}
{伽利略根据理想实验推论出，如果没有摩擦，在水平面上的物体，一旦具有某一个速度，将保持这个速度继续运动下去}

%% answer: A
%% memo:
\memoanswer{本题原卷及材料中未提供详解，待补充。}

\item
%% number: 2
%% paperName: 2006年广东高考第 2 题 4 分
%% typeId: 1
%% type: 单选题
%% chapter: 直线运动
%% point: 运动图象
%% method: 无
%% score: 4
%% degree: 550
%% duplicateId: 0
%% body:
a、b 两物体从同一位置沿同一直线运动，它们的速度图象如图所示，下列说法正确的是\xzanswer{C}

\onepicture[w=5.0cm]{figs/02.png}

\fourchoices[answer=C]
{a、b 加速时，物体 a 的加速度大于物体 b 的加速度}
{$20\Us$ 时，a、b 两物体相距最远}
{$60\Us$ 时，物体 a 在物体 b 的前方}
{$40\Us$ 时，a、b 两物体速度相等，相距 $200\Um$}

%% answer: C
%% memo:
\memoanswer{$v-t$ 图象中，图象的斜率表示加速度，图线和时间轴所夹的面积表示位移。a、b 物体加速时，由图象可知 b 物体的加速度大于 a 物体的加速度，故 A 错误；当两物体的速度相等时，距离最大，故 B 错误；$60\Us$ 时，物体 a 的图线与时间轴围成的面积大于物体 b 与时间轴围成的面积，所以物体 a 在物体 b 的前方，故 C 正确；$40\Us$ 时，a、b 两物体速度相等，a 物体的图线与时间轴围成的面积为 $s_{a}=40\times40-\frac{1}{2}\times30\times20=1300\Um$，b 物体的图线与时间轴围成的面积为 $s_{b}=\frac{1}{2}\times20\times40=400\Um$，相距 $s_{a}-s_{b}=1300\Um-400\Um=900\Um$，故 D 错误。}

\item
%% number: 3
%% paperName: 2006年广东高考第 3 题 4 分
%% typeId: 2
%% type: 多选题
%% chapter: 其他
%% point: 物理学史
%% method: 无
%% score: 4
%% degree: 850
%% duplicateId: 0
%% body:
下列说法正确的是\xzanswer{BCD}

\fourchoices[answer=BCD]
{康普顿发现了电子}
{卢瑟福提出了原子的核式结构模型}
{贝克勒尔发现了铀和含铀矿物的天然放射现象}
{伦琴发现了 X 射线}

%% answer: BCD
%% memo:
\memoanswer{康普顿发现了康普顿效应，汤姆孙发现了电子，故 A 错误；卢瑟福提出了原子的核式结构模型，故 B 正确；贝克勒尔发现了铀和含铀矿物的天然放射现象，伦琴发现了 X 射线，故 C、D 正确。}

\item
%% number: 4
%% paperName: 2006年广东高考第 4 题 4 分
%% typeId: 1
%% type: 单选题
%% chapter: 气体性质
%% point: 热力学第一定律
%% method: 无
%% score: 4
%% degree: 650
%% duplicateId: 0
%% body:
关于永动机和热力学定律的讨论，下列叙述正确的是\xzanswer{D}

\fourchoices[answer=D]
{第二类永动机违反能量守恒定律}
{如果物体从外界吸收了热量，则物体的内能一定增加}
{外界对物体做功，则物体的内能一定增加}
{做功和热传递都可以改变物体的内能，但从能量转化或转移的观点来看这两种改变方式是有区别的}

%% answer: D
%% memo:
\memoanswer{第二类永动机并不违反能量守恒定律，而是违反了热力学第二定律，跟热现象有关的宏观过程均具有方向性，故 A 错误；做功和热传递是改变物体内能的两种方式，做功是内能和其他形式的能之间的转化，热传递是内能之间的转移，内能是否变化取决于这两种方式，故 B、C 错误、D 正确。}

\item
%% number: 5
%% paperName: 2006年广东高考第 5 题 4 分
%% typeId: 2
%% type: 多选题
%% chapter: 原子和原子核
%% point: 核裂变
%% method: 无
%% score: 4
%% degree: 650
%% duplicateId: 0
%% body:
据新华社报道，由我国自行设计、研制的世界第一套全超导核聚变实验装置（又称“人造太阳”）已完成了首次工程调试。下列关于“人造太阳”的说法正确的是\xzanswer{AC}

\fourchoices[answer=AC]
{“人造太阳”的核反应方程是 $\ce{^{2}_{1}H + ^{3}_{1}H -> ^{4}_{2}He + ^{1}_{0}n}$}
{“人造太阳”的核反应方程是 $\ce{^{235}_{92}U + ^{1}_{0}n -> ^{141}_{56}Ba + ^{92}_{36}Kr + 3^{1}_{0}n}$}
{“人造太阳”释放的能量大小的计算公式是 $\Delta E=\Delta mc^{2}$}
{“人造太阳”核能大小的计算公式是 $E=\frac{1}{2}mc^{2}$}

%% answer: AC
%% memo:
\memoanswer{“人造太阳”是聚变反应，聚变反应是两个轻核结合成质量较大的原子核的核反应，故 A 正确、B 错误；释放的能量大小用爱因斯坦的质能方程 $\Delta E=\Delta mc^{2}$ 计算，故 C 正确、D 错误。}

\item
%% number: 6
%% paperName: 2006年广东高考第 6 题 4 分
%% typeId: 2
%% type: 多选题
%% chapter: 机械振动
%% point: 受迫振动共振
%% method: 无
%% score: 4
%% degree: 650
%% duplicateId: 0
%% body:
铺设铁轨时，每两根钢轨接缝处都必须留有一定的间隙，匀速运行的列车经过轨端接缝处时，车轮就会受到一次冲击。由于每一根钢轨长度相等，所以这个冲击力是周期性的，列车受到周期性的冲击做受迫振动。普通钢轨长为 $12.6\Um$，列车固有振动周期为 $0.315\Us$。下列说法正确的是\xzanswer{AD}

\fourchoices[answer=AD]
{列车的危险速率为 $40\Ums$}
{列车过桥需要减速，是为了防止列车发生共振现象}
{列车运行的振动频率和列车的固有频率总是相等的}
{增加钢轨的长度有利于列车高速运行}

%% answer: AD
%% memo:
\memoanswer{共振的条件是驱动力的频率等于系统的固有频率，由 $\frac{l}{v}=T$ 可求出危险车速为 $40\Ums$，故 A 正确；B、C 选项的说法过于绝对，故错误；增加钢轨的长度可以降低固有频率，有利于列车高速运行，故 D 正确。}

\item
%% number: 7
%% paperName: 2006年广东高考第 7 题 4 分
%% typeId: 1
%% type: 单选题
%% chapter: 光学
%% point: 光的折射
%% method: 无
%% score: 4
%% degree: 650
%% duplicateId: 0
%% body:
两束不同频率的单色光 a、b 从空气射入水中，发生了如图所示的折射现象（$\alpha>\beta$）。下列结论中正确的是\xzanswer{C}

\onepicture[w=5.0cm]{figs/07.png}

\fourchoices[answer=C]
{光束 b 的频率比光束 a 低}
{在水中的传播速度，光束 a 比光束 b 小}
{水对光束 a 的折射率比水对光束 b 的折射率小}
{若光束从水中射向空气，则光束 b 的临界角比光束 a 的临界角大}

%% answer: C
%% memo:
\memoanswer{由 $n=\frac{\sin\theta_{1}}{\sin\theta_{2}}$ 知，b 的折射率较大，则 b 的频率较大，在同种介质中传播速度较小，对同种介质的临界角较小，故 C 正确、A、B、D 错误。}

\item
%% number: 8
%% paperName: 2006年广东高考第 8 题 4 分
%% typeId: 2
%% type: 多选题
%% chapter: 气体性质
%% point: 热力学第一定律
%% method: 无
%% score: 4
%% degree: 650
%% duplicateId: 0
%% body:
如图所示为电冰箱的工作原理示意图。压缩机工作时，强迫致冷剂在冰箱内外的管道中不断循环。在蒸发器中致冷剂汽化吸收箱体内的热量，经过冷凝器时致冷剂液化，放出热量到箱体外。下列说法正确的是\xzanswer{BC}

\onepicture[w=7.5cm]{\includesvg{figs/08}}

\fourchoices[answer=BC]
{热量可以自发地从冰箱内传到冰箱外}
{电冰箱的致冷系统能够不断地把冰箱内的热量传到外界，是因为其消耗了电能}
{电冰箱的工作原理不违反热力学第一定律}
{电冰箱的工作原理违反热力学第一定律}

%% answer: BC
%% memo:
\memoanswer{由热力学第二定律知，热量不能自发地由低温物体传到高温物体，除非施加外部的影响和帮助。电冰箱把热量从低温的内部传到高温外部，需要压缩机的帮助并消耗电能，故 B、C 正确。}

\item
%% number: 9
%% paperName: 2006年广东高考第 9 题 4 分
%% typeId: 2
%% type: 多选题
%% chapter: 光学
%% point: 光的电磁说
%% method: 无
%% score: 4
%% degree: 650
%% duplicateId: 0
%% body:
目前雷达发射的电磁波频率多在 $200\UMHz$ 至 $1000\UMHz$ 的范围内。下列关于雷达和电磁波的说法正确的是\xzanswer{ACD}

\fourchoices[answer=ACD]
{真空中上述雷达发射的电磁波的波长范围在 $0.3\Um$ 至 $1.5\Um$ 之间}
{电磁波是由恒定不变的电场或磁场产生的}
{测出从发射电磁波到接收反射波的时间间隔可以确定雷达和目标的距离}
{波长越短的电磁波，反射性能越强}

%% answer: ACD
%% memo:
\memoanswer{据 $\lambda=\frac{c}{f}$，电磁波频率在 $200\UMHz$ 至 $1000\UMHz$ 范围内，则电磁波的波长范围在 $0.3\Um$ 至 $1.5\Um$ 之间，故 A 正确；电磁波的产生依据麦克斯韦的电磁场理论，变化的磁场产生电场，变化的电场产生磁场，故 B 错误；雷达是利用电磁波的反射原理，且波长越短，粒子性表现得越明显，反射性能越强，故 C、D 正确。}

\item
%% number: 10
%% paperName: 2006年广东高考第 10 题 4 分
%% typeId: 1
%% type: 单选题
%% chapter: 电磁感应
%% point: 楞次定律推广
%% method: 无
%% score: 4
%% degree: 450
%% duplicateId: 0
%% body:
如图所示，用一根长为 $L$、质量不计的细杆与一个上弧长为 $l_{0}$、下弧长为 $d_{0}$ 的金属线框的中点联结并悬挂于 O 点，悬点正下方存在一个上弧长为 $2l_{0}$、下弧长为 $2d_{0}$ 的方向垂直纸面向里的匀强磁场，且 $d_{0}\ll L$。先将线框拉开到如图所示位置，松手后让线框进入磁场，忽略空气阻力和摩擦。下列说法正确的是\xzanswer{D}

\onepicture[w=5.0cm]{figs/10.png}

\fourchoices[answer=D]
{金属线框进入磁场时感应电流的方向为 $a\to b\to c\to d\to a$}
{金属线框离开磁场时感应电流的方向为 $a\to d\to c\to b\to a$}
{金属线框 dc 边进入磁场与 ab 边离开磁场的速度大小总是相等}
{金属线框最终将在磁场内做简谐运动}

%% answer: D
%% memo:
\memoanswer{金属线框进入磁场时，由于电磁感应产生电流，由楞次定律判断电流的方向为 $a\to d\to c\to b\to a$，故 A 错误；金属线框离开磁场时由于电磁感应产生电流，由楞次定律判断电流的方向为 $a\to b\to c\to d\to a$，故 B 错误；由能量转化和守恒可知，金属线框 dc 边进入磁场与 ab 边离开磁场的速度大小不相等，故 C 错误；如此往复摆动，最终金属线框在匀强磁场内摆动，由于 $d_{0}\ll L$，单摆做简谐运动的条件是摆角很小，故最终在磁场内做简谐运动，D 正确。}

\item
%% number: 11
%% paperName: 2006年广东高考第 11 题 9 分
%% typeId: 4
%% type: 实验题
%% chapter: 曲线运动
%% point: 研究平抛运动实验
%% method: 无
%% score: 9
%% degree: 650
%% duplicateId: 0
%% body:
某同学设计了一个研究平抛运动的实验。实验装置示意图如图~\ref{2006广东11a} 所示，A 是一块平面木板，在其上等间隔地开凿出一组平行的插槽（图~\ref{2006广东11a} 中 $P_{0}P_{0}'$、$P_{1}P_{1}'$、$\ldots\ldots$），槽间距离均为 $d$。把覆盖复写纸的白纸铺贴在硬板 B 上。实验时依次将 B 板插入 A 板的各插槽中，每次让小球从斜轨的同一位置由静止释放。每打完一点后，把 B 板插入后一槽中并同时向纸面内侧平移距离 $d$。实验得到小球在白纸上打下的若干痕迹点，如图~\ref{2006广东11b} 所示。

\twopicture[labelA=2006广东11a, labelB=2006广东11b, w=7.5cm]
{figs/11a.png}
{figs/11b.png}

\begin{enumerate}
	\item
	实验前应对实验装置反复调节，直到 \tkanswer[3cm]{}。每次让小球从同一位置由静止释放，是为了 \tkanswer[4cm]{}。
	\item
	每次将 B 板向内侧平移距离 $d$，是为了 \tkanswer[5cm]{}。
	\item
	在图~\ref{2006广东11b} 中绘出小球做平抛运动的轨迹。
\end{enumerate}

%% answer:
\jdanswer{
\begin{enumerate}
	\item
	斜槽末端水平，保持小球水平抛出的初速度相同
	\item
	保持相邻痕迹点的水平距离大小相同，使轨迹连线与小球的平抛轨迹相同
	\item
	平抛运动的轨迹如图。
\end{enumerate}
}

%% memo:
\memoanswer{平抛运动可分解为水平方向的匀速直线运动和竖直方向的自由落体运动。实验前应对实验装置反复调节，直到斜槽末端水平；每次让小球从同一位置由静止释放，是为了保持小球水平抛出的初速度相同。每次将 B 板向内侧平移距离 $d$，是为了保持相邻痕迹点的水平距离大小相同，使轨迹连线与小球的平抛轨迹相同。该题考查了实验中的留迹法。

\onepicture[w=5.0cm]{figs/11_an.png}
}

\item
%% number: 12
%% paperName: 2006年广东高考第 12 题 11 分
%% typeId: 4
%% type: 实验题
%% chapter: 电路
%% point: 测电动势与内阻
%% method: 无
%% score: 11
%% degree: 450
%% duplicateId: 0
%% body:
某同学设计了一个如图~\ref{2006广东12a} 所示的实验电路，用以测定电源电动势和内阻，使用的实验器材为：待测干电池组（电动势约 $3\UV$）、电流表（量程 $0.6\UA$，内阻小于 $1\UO$）、电阻箱（$0\sim99.99\UO$）、滑动变阻器（$0\sim10\UO$）、单刀双掷开关、单刀单掷开关各一个及导线若干。考虑到干电池的内阻较小，电流表的内阻不能忽略。

\onepicture[num=a, showlabel=false, label=2006广东12a, w=4.5cm]{figs/12a.png}

\begin{enumerate}
	\item
	该同学按图~\ref{2006广东12a} 连线，通过控制开关状态，测得电流表内阻约为 $0.20\UO$。试分析该测量产生误差的原因是 \tkanswer[5cm]{}。
	\item
	简要写出利用图~\ref{2006广东12a} 所示电路测量电源电动势和内阻的实验步骤：
	①\tkanswer[6cm]{}；
	②\tkanswer[6cm]{}。
	\item
	图~\ref{2006广东12b} 是由实验数据绘出的 $\frac{1}{I}-R$ 图象，由此求出待测干电池组的电动势 $E=$ \tkanswer[1.5cm]{}$\UV$、内阻 $r=$ \tkanswer[1.5cm]{}$\UO$。（计算结果保留三位有效数字）
\end{enumerate}

\onepicture[num=b, showlabel=false, label=2006广东12b, w=6.5cm, align=right]{figs/12b.png}

%% answer:
\jdanswer{
\begin{enumerate}
	\item
	并联电阻箱后线路总阻值减小，从而造成总电流增大
	\item
	①调节电阻箱 $R$，断开开关 K，将开关 S 接 D，记录电阻箱的阻值和电流表示数；②断开开关 D，再次调节电阻箱 $R$，将开关 S 接 D，记录电阻箱的阻值和电流表示数
	\item
	$2.81$，$2.33$
\end{enumerate}
}

%% memo:
\memoanswer{由测定电源电动势和内阻实验的原理知，此种接法出现误差的原因是电流表的分压作用。而 $R=r+r_{\text{安}}$，$\frac{1}{I}-R$ 图线的斜率表示电源电动势的倒数，据此得出电动势 $E=2.81\UV$，内阻 $r=2.33\UO$。}

\item
%% number: 13
%% paperName: 2006年广东高考第 13 题 7 分
%% typeId: 6
%% type: 计算题
%% chapter: 功和能
%% point: 动能定理
%% method: 无
%% score: 7
%% degree: 650
%% duplicateId: 0
%% body:
\begin{enumerate}
	\item
	人们发现光电效应具有瞬时性和对各种金属都存在极限频率的规律。请问谁提出了何种学说很好地解释了上述规律？已知锌的逸出功为 $3.34\UeV$，用某单色紫外线照射锌板时，逸出光电子的最大速度为 $10^{6}\Ums$，求该紫外线的波长 $\lambda$（电子质量 $m_{e}=9.11\times10^{-31}\Ukg$，普朗克常量 $h=6.63\times10^{-34}\UJcdots$，$1\UeV=1.6\times10^{-19}\UJ$）。
	\item
	风力发电是一种环保的电能获取方式。图为某风力发电站外观图。设计每台风力发电机的功率为 $40\UkW$。实验测得风的动能转化为电能的效率约为 $20\%$，空气的密度是 $1.29\Ukgmc$，当地水平风速约为 $10\Ums$，问风力发电机的叶片长度约为多少才能满足设计要求？

	\onepicture[w=5.0cm, align=right]{figs/13.jpg}
\end{enumerate}

%% answer:
\jdanswer{
\begin{enumerate}
	\item
	爱因斯坦；$3.23\times10^{-7}\Um$
	\item
	$10\Um$
\end{enumerate}
}

%% memo:
\memoanswer{
（1）爱因斯坦提出的光子说很好地解释了光电效应现象。由光电效应方程得
\[
\frac{hc}{\lambda}=W_{0}+\frac{1}{2}m_{e}v^{2},
\]
解得 $\lambda=3.23\times10^{-7}\Um$。

（2）风的动能为 $E_{\mathrm{k}}=\frac{1}{2}mv^{2}$ ①，$t$ 时间内风通过以叶片为半径的圆的质量为 $m=\rho V=\rho\pi l^{2}vt$ ②，则发电机的功率为 $P=\frac{\eta E_{\mathrm{k}}}{t}=\frac{1}{2}\rho\pi l^{2}v^{3}\eta$ ③，由①②③得 $l\approx10\Um$ ④。
}

\item
%% number: 14
%% paperName: 2006年广东高考第 14 题 12 分
%% typeId: 6
%% type: 计算题
%% chapter: 电磁感应
%% point: 远距离输电
%% method: 无
%% score: 12
%% degree: 650
%% duplicateId: 0
%% body:
某发电站的输出功率为 $10^{4}\UkW$，输出电压为 $4\UkV$，通过理想变压器升压后向 $80\Ukm$ 远处供电。已知输电导线的电阻率为 $\rho=2.4\times10^{-8}\UOm$，导线横截面积为 $1.5\times10^{-4}\Umq$，输电线路损失的功率为输出功率的 $4\%$，求：

\begin{enumerate}
	\item
	升压变压器的输出电压；
	\item
	输电线路上的电压损失。
\end{enumerate}

%% answer:
\jdanswer{
\begin{enumerate}
	\item
	$8\times10^{4}\UV$
	\item
	$3200\UV$
\end{enumerate}
}

%% memo:
\memoanswer{
（1）由电阻定律知，输电线的电阻为
\[
R=\rho\frac{2l}{S}\quad\text{①},
\]
输电线损失的功率为 $P_{\text{损}}=P_{1}\times4\%$ ②，又 $P_{\text{损}}=I_{2}^{2}R$ ③。由理想变压器的关系知 $P_{2}=P_{1}$ ④，又 $P_{2}=I_{2}U_{2}$ ⑤，由①②③④⑤得升压变压器的输出电压为 $U_{2}=8\times10^{4}\UV$ ⑥。

（2）输电线路上的电压损失为 $U_{\text{损}}=I_{2}R=3200\UV$ ⑦。
}

\item
%% number: 15
%% paperName: 2006年广东高考第 15 题 14 分
%% typeId: 6
%% type: 计算题
%% chapter: 运动和力
%% point: 牛顿定律的应用
%% method: 无
%% score: 14
%% degree: 450
%% duplicateId: 0
%% body:
一个质量为 $4\Ukg$ 的物体静止在足够大的水平地面上，物体与地面间的动摩擦因数 $\mu=0.1$。从 $t=0$ 开始，物体受到一个大小和方向呈周期性变化的水平力 $F$ 作用，力 $F$ 随时间的变化规律如图所示。求 83 秒内物体的位移大小和力 $F$ 对物体所做的功。

\onepicture[w=8.0cm, align=right]{figs/15.png}

%% answer:
\jdanswer{$x=167\Um$，$W=681\UJ$}

%% memo:
\memoanswer{
当物体在前半周期时，由牛顿第二定律得
\[
F_{1}-\mu mg=ma_{1},\qquad a_{1}=\frac{F_{1}-\mu mg}{m}=\frac{12-0.1\times4\times10}{4}\Umsq=2\Umsq;
\]
当物体在后半周期时，由牛顿第二定律得
\[
F_{2}+\mu mg=ma_{2},\qquad a_{2}=\frac{F_{2}+\mu mg}{m}=\frac{4+0.1\times4\times10}{4}\Umsq=2\Umsq.
\]
前半周期和后半周期位移相等，$x_{1}=\frac{1}{2}a_{1}t^{2}=0.5\times2\times2^{2}\Um=4\Um$，一个周期的位移为 $8\Um$，最后 $1\Us$ 的位移为 $3\Um$。83 秒内物体的位移大小为
\[
x=20\times8\Um+4\Um+3\Um=167\Um.
\]
一个周期 $F$ 做的功为 $W_{1}=(F_{1}-F_{2})x_{1}=(12-4)\times4\UJ=32\UJ$，力 $F$ 对物体所做的功
\[
W=20\times32\UJ+12\times4\UJ-4\times3\UJ=681\UJ.
\]
}

\item
%% number: 16
%% paperName: 2006年广东高考第 16 题 16 分
%% typeId: 6
%% type: 计算题
%% chapter: 电磁感应
%% point: 电磁感应能量分析
%% method: 无
%% score: 16
%% degree: 450
%% duplicateId: 0
%% body:
如图所示，在磁感应强度大小为 $B$、方向垂直向上的匀强磁场中，有一上、下两层均与水平面平行的“U”型光滑金属导轨，在导轨面上各放一根完全相同的质量为 $m$ 的匀质金属杆 $\mathrm{A}_{1}$ 和 $\mathrm{A}_{2}$，开始时两根金属杆位于同一竖直面内且杆与轨道垂直。设两导轨面相距为 $H$，导轨宽为 $L$，导轨足够长且电阻不计，金属杆单位长度的电阻为 $r$。现有一质量为 $\frac{m}{2}$ 的不带电小球以水平向右的速度 $v_{0}$ 撞击杆 $\mathrm{A}_{1}$ 的中点，撞击后小球反弹落到下层面上的 C 点。C 点与杆 $\mathrm{A}_{2}$ 初始位置相距为 $s$。求：

\begin{enumerate}
	\item
	回路内感应电流的最大值；
	\item
	整个运动过程中感应电流最多产生了多少热量；
	\item
	当杆 $\mathrm{A}_{2}$ 与杆 $\mathrm{A}_{1}$ 的速度比为 $1:3$ 时，$\mathrm{A}_{2}$ 受到的安培力大小。
\end{enumerate}

\onepicture[w=7.5cm, align=right]{figs/16.png}

%% answer:
\jdanswer{
\begin{enumerate}
	\item
	$I=\frac{B\left(v_{0}+s\sqrt{\frac{g}{2H}}\right)}{4r}$
	\item
	$Q=\frac{1}{16}m\left(v_{0}+s\sqrt{\frac{g}{2H}}\right)^{2}$
	\item
	$F_{2}=\frac{B^{2}L\left(v_{0}+s\sqrt{\frac{g}{2H}}\right)}{8r}$
\end{enumerate}
}

%% memo:
\memoanswer{
（1）小球撞击杆瞬间动量守恒，之后做平抛运动。设小球碰撞后速度大小为 $v_{1}$，杆获得速度大小为 $v_{2}$，则
\[
\frac{m}{2}v_{0}=-\frac{m}{2}v_{1}+mv_{2},\qquad s=v_{1}t,\qquad H=\frac{1}{2}gt^{2},
\]
联立解得 $v_{2}=\frac{1}{2}\left(v_{0}+s\sqrt{\frac{g}{2H}}\right)$。杆在磁场中运动，其最大的电动势为 $E_{1}=BLv_{2}$，最大电流为 $I_{\max}=\frac{E_{1}}{2Lr}$，联立解得 $I_{\max}=\frac{B\left(v_{0}+s\sqrt{\frac{g}{2H}}\right)}{4r}$。

（2）两金属杆在磁场中运动始终满足动量守恒，两杆最终速度相同设为 $v'$，则有 $mv_{2}=2mv'$，由功能关系得 $Q=\frac{1}{2}mv_{2}^{2}-\frac{1}{2}\cdot2mv'^{2}$，联立解得 $Q=\frac{1}{16}m\left(v_{0}+s\sqrt{\frac{g}{2H}}\right)^{2}$。

（3）设杆 $\mathrm{A}_{2}$ 和 $\mathrm{A}_{1}$ 的速度大小分别为 $v$ 和 $3v$，由动量守恒得 $mv_{2}=mv+m\cdot3v$。由法拉第电磁感应定律得 $E_{2}=BL(3v-v)$，又 $I=\frac{E_{2}}{2Lr}$，则安培力的大小为 $F=BIL$，联立解得 $F=\frac{B^{2}L\left(v_{0}+s\sqrt{\frac{g}{2H}}\right)}{8r}$。
}

\item
%% number: 17
%% paperName: 2006年广东高考第 17 题 16 分
%% typeId: 6
%% type: 计算题
%% chapter: 曲线运动
%% point: 人造地球卫星
%% method: 无
%% score: 16
%% degree: 450
%% duplicateId: 0
%% body:
宇宙中存在一些离其他恒星较远的、由质量相等的三颗星组成的三星系统，通常可忽略其他星体对它们的引力作用。已观测到稳定的三星系统存在两种基本的构成形式：一种是三颗星位于同一直线上，两颗星围绕中央星在同一半径为 $R$ 的圆轨道上运行；另一种形式是三颗星位于等边三角形的三个顶点上，并沿外接于等边三角形的圆形轨道运行。设每个星体的质量均为 $m$。

\begin{enumerate}
	\item
	试求第一种形式下，星体运动的线速度和周期。
	\item
	假设两种形式星体的运动周期相同，第二种形式下星体之间的距离应为多少？
\end{enumerate}

%% answer:
\jdanswer{
\begin{enumerate}
	\item
	$v=\sqrt{\frac{5Gm}{4R}}$，$T=4\pi R\sqrt{\frac{R}{5Gm}}$
	\item
	$l=\left(\frac{12}{5}\right)^{\frac{1}{3}}R$
\end{enumerate}
}

%% memo:
\memoanswer{
（1）对于第一种运动情况，以某个运动星体为研究对象，由牛顿第二定律和万有引力定律得
\[
F_{1}=\frac{Gm^{2}}{R^{2}},\qquad F_{2}=\frac{Gm^{2}}{(2R)^{2}},
\]
则 $F_{1}+F_{2}=\frac{mv^{2}}{R}$ ①，运动星体的线速度大小为 $v=\frac{\sqrt{5GmR}}{2R}$ ②。设星体运动的周期为 $T$，有 $T=\frac{2\pi R}{v}$ ③，联立解得 $T=4\pi\sqrt{\frac{R^{3}}{5Gm}}$ ④。

\onepicture[w=3.5cm]{figs/17_an1.png}

（2）设第二种形式星体之间的距离为 $r$，则三个星体做圆周运动的半径为 $R'=\frac{r/2}{\cos30^{\circ}}$ ⑤。由于星体做圆周运动所需要的向心力靠其他两个星体的万有引力的合力提供，由力的合成和牛顿运动定律有 $F_{\text{合}}=2\frac{Gm^{2}}{r^{2}}\cos30^{\circ}$ ⑥，$F_{\text{合}}=m\frac{4\pi^{2}}{T^{2}}R'$ ⑦，由④⑤⑥⑦式得 $l=\left(\frac{12}{5}\right)^{\frac{1}{3}}R$ ⑧。

\onepicture[w=3.5cm]{figs/17_an2.png}
}

\item
%% number: 18
%% paperName: 2006年广东高考第 18 题 17 分
%% typeId: 6
%% type: 计算题
%% chapter: 磁场
%% point: 洛伦兹力
%% method: 无
%% score: 17
%% degree: 450
%% duplicateId: 0
%% body:
在光滑绝缘的水平桌面上，有两个质量均为 $m$、电荷量为 $+q$ 的完全相同的带电粒子 $\mathrm{P}_{1}$ 和 $\mathrm{P}_{2}$，在小孔 A 处以初速度为零先后释放。在平行板间距为 $d$ 的匀强电场中加速后，$\mathrm{P}_{1}$ 从 C 处对着圆心进入半径为 $R$ 的固定圆筒中（筒壁上的小孔 C 只能容一个粒子通过），圆筒内有垂直水平面向上的磁感应强度为 $B$ 的匀强磁场。$\mathrm{P}_{1}$ 每次与筒壁发生碰撞均无电荷迁移，$\mathrm{P}_{1}$ 进入磁场第一次与筒壁碰撞点为 D，$\angle COD=\theta$，如图所示。延后释放的 $\mathrm{P}_{2}$，将第一次欲逃逸出圆筒的 $\mathrm{P}_{1}$ 正碰圆筒内，此次碰撞刚结束，立即改变平行板间的电压，并利用 $\mathrm{P}_{2}$ 与 $\mathrm{P}_{1}$ 之后的碰撞，将 $\mathrm{P}_{1}$ 限制在圆筒内运动。碰撞过程均无机械能损失。设 $d=\frac{5}{8}\pi R$，求：在 $\mathrm{P}_{2}$ 和 $\mathrm{P}_{1}$ 相邻两次碰撞时间间隔内，粒子 $\mathrm{P}_{1}$ 与筒壁的可能碰撞次数。

\onepicture[w=7.5cm, align=right]{figs/18.png}

附：部分三角函数值

\begin{table}[htbp]
\centering
\begin{tblr}{|c|c|c|c|c|c|c|c|c|c|}
\hline
$\varphi$ & $\frac{2\pi}{5}$ & $\frac{\pi}{3}$ & $\frac{\pi}{4}$ & $\frac{\pi}{5}$ & $\frac{\pi}{6}$ & $\frac{\pi}{7}$ & $\frac{\pi}{8}$ & $\frac{\pi}{9}$ & $\frac{\pi}{10}$ \\
\hline
$\tan\varphi$ & 3.08 & 1.73 & 1.00 & 0.73 & 0.58 & 0.48 & 0.41 & 0.36 & 0.32 \\
\hline
\end{tblr}
\end{table}

%% answer:
\jdanswer{粒子与圆筒碰撞的可能次数为 2，3，4，5，6，7。}

%% memo:
\memoanswer{
设开始加速电压为 $U_{1}$，$\mathrm{P}_{1}$、$\mathrm{P}_{2}$ 第一次碰撞后电压变为 $U_{2}$。由动能定理得 $qU_{1}=\frac{1}{2}mv_{0}^{2}$ ①，洛伦兹力提供向心力得 $qv_{0}B=m\frac{v_{0}^{2}}{r}$ ②。设粒子在圆筒中与筒壁碰撞 $(n-1)$ 次，转了 $K$ 圈后从 C 点射出，粒子运动时间为 $t_{1}$，由几何关系得 $n\theta=2K\pi$ ③，$\tan\frac{\theta}{2}=\frac{r}{R}$ ④，粒子在磁场中运动的周期为 $T_{0}=\frac{2\pi r}{v_{0}}$，联立解得 $t_{1}=n\times\frac{\pi-\theta}{2\pi}T_{0}$ ⑤。

\onepicture[w=5.0cm]{figs/18_an.png}

$\mathrm{P}_{2}$ 与 $\mathrm{P}_{1}$ 发生弹性碰撞，由动量与能量守恒得 $mv_{0}+m(-v_{0})=mv_{1}+mv_{2}$，$\frac{1}{2}mv_{0}^{2}+\frac{1}{2}mv_{0}^{2}=\frac{1}{2}mv_{1}^{2}+\frac{1}{2}mv_{2}^{2}$，联立解得 $v_{1}=-v_{2}$ ⑥，$v_{1}=-v_{0}$、$v_{2}=v_{0}$ ⑦，即两粒子速度发生交换。

$\mathrm{P}_{2}$ 在电场中运动时间为 $t_{2}$，$\mathrm{P}_{1}$、$\mathrm{P}_{2}$ 碰撞后速度发生交换后有 $t_{2}=\frac{2v_{0}}{a}=\frac{2mdv_{0}}{qU_{2}}$ ⑧。欲使 $\mathrm{P}_{2}$ 与 $\mathrm{P}_{1}$ 碰撞后不穿出 A 点，与 $\mathrm{P}_{1}$ 碰撞后 $\mathrm{P}_{2}$ 返回距离为 $s$，有 $\frac{U_{2}}{U_{1}}=\frac{d}{s}\geq1$，得 $U_{2}\geq U_{1}$ ⑨；欲使 $\mathrm{P}_{1}$ 不逃逸出磁场有 $t_{1}=t_{2}$ ⑩。由⑤⑧⑩得 $U_{2}=\frac{2dBv_{0}}{\pi(n-2K)}$（11），由①②④⑨（11）得 $\frac{U_{2}}{U_{1}}=\frac{d}{R}\times\frac{4}{\pi(n-2K)\tan\frac{K\pi}{n}}\geq1$（12）。

通过讨论：只有当 $K=1$、$n=3,4,5,6,7,8$ 时不等式成立，所以粒子与圆筒碰撞的可能次数为 2，3，4，5，6，7。
}

\end{enumerate}

\end{document}
